% preamble-exam-zh.tex — 共享 preamble
% 用于所有 exam-zh 模板
%
% 样式策略：
%   屏幕 → 语义彩色（蓝=关键想法，红=易错，绿=路线，灰=提示/教师备注）
%   打印 → 自动降级为黑白（通过 print mode 或 monochrome 选项）

\documentclass{exam-zh}

% ---------- 基础配置 ----------
\examsetup{
  page/size = a4paper,
  page/show-head = false,
  page/show-foot = false,
  sealline/show = false,
  font = times,
  math-font = xits,
}

\usepackage{tcolorbox}
\usepackage{needspace}
\tcbuselibrary{skins,breakable}

% ---------- 打印/屏幕颜色切换 ----------
% 用法：\documentclass[monochrome]{exam-zh} 即可切换为纯黑白
% 注意：不要使用 \if@xxx 放在 makeatother 外部；这里改成普通条件名。
\newif\ifeduprintcolor
\eduprintcolortrue

\makeatletter
\@ifclasswith{exam-zh}{monochrome}{\eduprintcolorfalse}{}
\makeatother

% ---------- 语义颜色定义 ----------
% 屏幕：有色彩区分；打印：降级为灰度
\ifeduprintcolor
  % --- 屏幕彩色模式 ---
  \definecolor{edu-blue}{RGB}{47, 82, 122}       % 关键想法
  \definecolor{edu-blue-bg}{RGB}{243, 246, 252}
  \definecolor{edu-red}{RGB}{180, 40, 40}         % 易错提醒
  \definecolor{edu-red-bg}{RGB}{255, 245, 240}
  \definecolor{edu-green}{RGB}{34, 100, 60}       % 路线图
  \definecolor{edu-green-bg}{RGB}{242, 250, 245}
  \definecolor{edu-orange}{RGB}{160, 90, 0}       % 训练目标
  \definecolor{edu-orange-bg}{RGB}{255, 248, 235}
  \definecolor{edu-gray}{RGB}{100, 100, 100}      % 提示/教师备注
  \definecolor{edu-gray-bg}{RGB}{248, 248, 248}
  \definecolor{edu-purple}{RGB}{90, 50, 120}      % 思考题
  \definecolor{edu-purple-bg}{RGB}{248, 244, 252}
\else
  % --- 打印黑白模式 ---
  \definecolor{edu-blue}{gray}{0.15}
  \definecolor{edu-blue-bg}{gray}{0.96}
  \definecolor{edu-red}{gray}{0.25}
  \definecolor{edu-red-bg}{gray}{0.97}
  \definecolor{edu-green}{gray}{0.20}
  \definecolor{edu-green-bg}{gray}{0.97}
  \definecolor{edu-orange}{gray}{0.25}
  \definecolor{edu-orange-bg}{gray}{0.98}
  \definecolor{edu-gray}{gray}{0.40}
  \definecolor{edu-gray-bg}{gray}{0.97}
  \definecolor{edu-purple}{gray}{0.30}
  \definecolor{edu-purple-bg}{gray}{0.98}
\fi

% ---------- 自定义教学环境 ----------
% 共性：enhanced + breakable（跨页不断裂）

% 原题卡片 — 强边框，突出显示
\newtcolorbox{problemcard}{
  enhanced, breakable,
  colback=edu-blue-bg,
  colframe=edu-blue,
  boxrule=1.2pt,
  arc=0pt,
  left=3mm, right=3mm, top=3mm, bottom=3mm,
}

% 解题路线图 — 左边线 + 浅底
\newtcolorbox{routemap}{
  enhanced, breakable,
  colback=edu-green-bg,
  colframe=edu-green!30,
  boxrule=0pt,
  borderline west={2.5pt}{0pt}{edu-green},
  left=3mm, right=3mm, top=3mm, bottom=3mm,
}

% 关键想法 — 蓝色左边线
\newtcolorbox{keyideabox}{
  enhanced, breakable,
  colback=edu-blue-bg,
  colframe=edu-blue!30,
  boxrule=0pt,
  borderline west={2.5pt}{0pt}{edu-blue},
  left=3mm, right=3mm, top=3mm, bottom=3mm,
}

% 易错提醒 — 红色左边线
\newtcolorbox{mistakebox}{
  enhanced, breakable,
  colback=edu-red-bg,
  colframe=edu-red!20,
  boxrule=0pt,
  borderline west={3pt}{0pt}{edu-red},
  left=3mm, right=3mm, top=2.5mm, bottom=2.5mm,
}

% 提示 — 灰色虚线左边线
\newtcolorbox{hintbox}{
  enhanced, breakable,
  colback=edu-gray-bg,
  colframe=edu-gray!20,
  boxrule=0pt,
  borderline west={2pt}{0pt}{edu-gray, dashed},
  left=3mm, right=3mm, top=2mm, bottom=2mm,
}

% 思考题 — 紫色虚线左边线
\newtcolorbox{thinkbox}{
  enhanced, breakable,
  colback=edu-purple-bg,
  colframe=edu-purple!20,
  boxrule=0pt,
  borderline west={2pt}{0pt}{edu-purple, dashed},
  left=2.5mm, right=2.5mm, top=2.5mm, bottom=2.5mm,
}

% 教师备注 — 灰色左边线，小字
\newtcolorbox{teachernote}{
  enhanced, breakable,
  colback=edu-gray-bg,
  colframe=edu-gray!15,
  boxrule=0pt,
  borderline west={2pt}{0pt}{edu-gray!60},
  left=2.5mm, right=2.5mm, top=2.5mm, bottom=2.5mm,
  fontupper=\small,
  colupper=black!60,
}

% 训练目标 — 橙色左边线，小字
\newtcolorbox{traininggoal}{
  enhanced, breakable,
  colback=edu-orange-bg,
  colframe=edu-orange!15,
  boxrule=0pt,
  borderline west={1.5pt}{0pt}{edu-orange},
  left=2mm, right=2mm, top=1mm, bottom=1mm,
  fontupper=\small,
}

% 答题区 — 无边框，纯留白
\newcommand{\answerarea}[1][40mm]{%
  \vspace{2mm}%
  \begin{tcolorbox}[
    enhanced, breakable,
    colback=white,
    colframe=white,
    boxrule=0pt,
    height=#1,
    valign=top,
    bottom=0mm,
    after skip=0mm,
  ]
  \end{tcolorbox}%
}

% ---------- 字体设置 ----------
% exam-zh (ctexbook) already sets CJK fonts via ctex.
% Only override if specific fonts are needed.
% macOS: Songti SC, STHeiti, STFangsong
% Linux: SimSun, SimHei, FangSong

% ---------- 页面设置 ----------
% exam-zh (ctexbook) already loads geometry.
% Override margins if needed:
% \geometry{top=16mm, bottom=16mm, left=14mm, right=14mm}

\examsetup{
  question/show-answer = true,
  fillin/show-answer = true,
  solution/show-solution = show-stay,
}

% ---------- 教师版解答样式 ----------
\newtcolorbox{solutionblock}{
  enhanced, breakable,
  colback=edu-blue-bg,
  colframe=edu-blue!25,
  boxrule=0pt,
  borderline west={2pt}{0pt}{edu-blue!50},
  arc=0pt,
  left=3mm, right=3mm, top=2mm, bottom=2mm,
  before skip=2mm,
  after skip=3mm,
  fontupper=\small,
}

% 解答步骤编号
\newcounter{solstep}
\newcommand{\solstep}[2]{%
  \stepcounter{solstep}%
  \par\medskip
  \noindent\hangindent=6mm
  {\color{edu-blue}\bfseries\textcircled{\scriptsize\thesolstep}\enspace #1}\enspace #2\par
}

\begin{document}

\section*{三垂直全等——专题练习（教师版）}

\section*{一、选择题（每题 3 分，共 12 分）}
（直接给出直角顶点和一个锐角顶点坐标，用一线三垂直求第三顶点）


\begin{question}[points=3]
  $A(2, 0)$，$B(0, 1)$，等腰 Rt$\triangle ABC$，$\angle A = 90°$，$C$ 在第一象限，则 $C$ 的坐标为
  \begin{choices}
    \item $(3, 2)$
    \item $(2, 3)$
    \item $(3, 3)$
    \item $(1, 2)$
  \end{choices}
  \paren[A]
  \begin{solutionblock}
    过直角顶点 $A$ 作水平线，过 $B$、$C$ 分别作这条水平线的垂线。一线三垂直全等后，$AB$ 的水平差 $2$ 与竖直差 $1$ 互换到 $AC$ 上。
    $C$ 在第一象限，所以从 $A(2,0)$ 向右 $1$、向上 $2$，得 $C(3,2)$。
  \end{solutionblock}
\end{question}



\begin{question}[points=3]
  $A(3, 0)$，$B(1, 3)$，等腰 Rt$\triangle ABC$，$\angle B = 90°$，$C$ 在第一象限，则 $C$ 的坐标为
  \begin{choices}
    \item $(6, 2)$
    \item $(2, 6)$
    \item $(6, 6)$
    \item $(4, 5)$
  \end{choices}
  \paren[D]
  \begin{solutionblock}
    直角顶点是 $B$。从 $B(1,3)$ 到 $A(3,0)$，水平差 $2$，竖直差 $3$。一线三垂直全等后，另一个直角边从 $B$ 出发应交换成“水平 $3$、竖直 $2$”。
    要在第一象限并与图形方向相符，取向右 $3$、向上 $2$，所以 $C(4,5)$。
  \end{solutionblock}
\end{question}



\begin{question}[points=3]
  已知$A(4, 0)$，$B(2, 8)$，等腰 Rt$\triangle ABC$，$\angle C = 90°$，$C$ 在第一象限，则 $C$ 的坐标为
  \begin{choices}
    \item $(7, 5)$
    \item $(4, 7)$
    \item $(7, 7)$
    \item $(3, 4)$
  \end{choices}
  \paren[A]
  \begin{solutionblock}
    $\angle C=90°$ 时，$AB$ 是等腰直角三角形的斜边。斜边中点为 $M(3,4)$，且 $CM=AM=BM$。
    $A$ 到 $M$ 的水平差 $1$、竖直差 $4$，过 $M$ 作与 $AB$ 垂直的方向后，$C$ 相对 $M$ 的差为水平 $4$、竖直 $1$。
    两个候选点为 $(7,5)$ 和 $(-1,3)$，题目要求第一象限，取 $C(7,5)$。
  \end{solutionblock}
\end{question}



\begin{question}[points=3]
  $A(1, 0)$，$B(0, 2)$，等腰 Rt$\triangle ABC$，在坐标平面内，有几个$C$点?
  \begin{choices}
    \item 2个
    \item 4个
    \item 6个
    \item 8个
  \end{choices}
  \paren[C]
  \begin{solutionblock}
    分直角位置讨论。若直角在 $A$，以 $AB$ 为一条直角边，可以在 $AB$ 两侧各作一个点 $C$，共 $2$ 个；若直角在 $B$，同理共 $2$ 个；若直角在 $C$，则 $AB$ 是斜边，斜边两侧各有一个等腰直角顶点，共 $2$ 个。
    合计 $2+2+2=6$ 个。
  \end{solutionblock}
\end{question}


\section*{二、填空题（每题 4 分，共 24 分）}


\begin{question}[points=4]
  $A(2, 1)$，$B(0, 3)$，等腰 Rt$\triangle ABC$，$\angle A = 90°$，$C$ 在第一象限，则 $C$ 的坐标为\fillin。
  \paren[$(4,\,3)$]
  \begin{solutionblock}
    直角顶点 $A$，从 $A$ 到 $B$ 的水平差 $2$、竖直差 $2$。一线三垂直全等后，另一条直角边也对应 $2$ 和 $2$。
    取第一象限且在 $A$ 右上方的点，$C(4,3)$。
  \end{solutionblock}
\end{question}



\begin{question}[points=4]
  直线 $y = -2x + 6$ 与 $x$ 轴、$y$ 轴分别交于 $A$、$B$，等腰 Rt$\triangle ABC$ 中 $\angle BAC = 90°$，$C$ 在第一象限，则 $C$ 的坐标为\fillin。
  \paren[$(9,\,3)$]
  \begin{solutionblock}
    交点为 $A(3,0)$，$B(0,6)$。从 $A$ 到 $B$ 的水平差 $3$、竖直差 $6$。
    过 $A$ 作水平辅助线，过 $B$、$C$ 作垂线，一线三垂直全等后，$AC$ 的水平差为 $6$，竖直差为 $3$。
    $C$ 在第一象限，故 $C(3+6,\,0+3)=(9,3)$。
  \end{solutionblock}
\end{question}



\begin{question}[points=4]
  已知 $A(1,2)$，$C(6,5)$。以 $AC$ 为直角边作等腰 Rt$\triangle ACE$，$\angle CAE = 90°$，且 $E$ 在 $A$ 的左上方，则 $E$ 的坐标为\fillin。
  \paren[$(-2,\,7)$]
  \begin{solutionblock}
    从 $A(1,2)$ 到 $C(6,5)$，水平差为 $5$，竖直差为 $3$。以 $A$ 为直角顶点作一线三垂直，另一条直角边的水平、竖直差互换为 $3$ 和 $5$。
    $E$ 在 $A$ 的左上方，所以从 $A$ 向左 $3$、向上 $5$，得 $E(-2,7)$。
  \end{solutionblock}
\end{question}



\begin{question}[points=4]
  已知 $B(7,1)$，$D(2,4)$。以 $BD$ 为直角边作等腰 Rt$\triangle BDE$，$\angle DBE = 90°$，且 $E$ 在 $B$ 的右上方，则 $E$ 的坐标为\fillin。
  \paren[$(10,\,6)$]
  \begin{solutionblock}
    从 $B(7,1)$ 到 $D(2,4)$，水平差为 $5$，竖直差为 $3$。一线三垂直全等后，$BE$ 的水平、竖直差互换为 $3$ 和 $5$。
    $E$ 在 $B$ 的右上方，所以 $E(7+3,1+5)=(10,6)$。
  \end{solutionblock}
\end{question}



\begin{question}[points=4]
  $A(1, 0)$，$B(4, 4)$，等腰 Rt$\triangle ABD$ 中 $\angle ABD = 90°$，$D$ 在第一象限，则 $D$ 的坐标为\fillin。
  \paren[$(8,\,1)$]
  \begin{solutionblock}
    直角顶点是 $B$。从 $B(4,4)$ 到 $A(1,0)$，水平差 $3$、竖直差 $4$。一线三垂直全等后，$BD$ 的水平差应为 $4$，竖直差应为 $3$。
    $D$ 在第一象限，取从 $B$ 向右 $4$、向下 $3$，得 $D(8,1)$。
  \end{solutionblock}
\end{question}



\begin{question}[points=4]
  等腰 Rt$\triangle ABC$ 中 $\angle C = 90°$，$AC = BC = 2$，$P$ 为 $\triangle ABC$ 内部一点，$AP = BP = CP = \sqrt{2}$，则 $\angle APC = $\fillin$°$。
  \paren[90]
  \begin{solutionblock}
    因为 $AP=BP=CP$，点 $P$ 到三个顶点距离相等，是这个直角三角形斜边 $AB$ 的中点。
    取最简单的图形：让 $C$ 为直角顶点，$AC=BC=2$，则 $P$ 是斜边中点，$PA=PC=\sqrt2$。此时 $\triangle APC$ 也是等腰直角三角形，所以 $\angle APC=90°$。
  \end{solutionblock}
\end{question}


\section*{三、解答题（每题 10 分，共 20 分）}

\needspace{8\baselineskip}

\begin{problem}[points=10]
  直线 $y = -\dfrac{5}{3}x + 10$ 与 $x$ 轴、$y$ 轴分别交于 $A$、$B$。以 $AB$ 为直角边在第一象限作等腰 Rt$\triangle ABC$，$\angle BAC = 90°$。\\
(1) 求点 $C$ 的坐标；\\
(2) 求 $\triangle ABC$ 的面积。

  \answerarea[70mm]
  \setcounter{solstep}{0}
  \begin{solutionblock}
    \textbf{答案：}$C(16,\,6)$，$S_{\triangle ABC}=68$\par\medskip
    \solstep{求交点}{令 $y=0$ 得 $A(6,0)$；令 $x=0$ 得 $B(0,10)$。}
    \solstep{作一线三垂直}{过直角顶点 $A$ 作水平辅助线，过 $B$、$C$ 分别作垂线。$AB$ 的水平差为 $6$，竖直差为 $10$，全等后 $AC$ 对应水平差 $10$、竖直差 $6$。}
    \solstep{求坐标}{$C$ 在第一象限，所以从 $A(6,0)$ 向右 $10$、向上 $6$，得 $C(16,6)$。}
    \solstep{求面积}{$AB=AC$，且 $\angle BAC=90°$。又 $AB^2=6^2+10^2=136$，所以 $S_{\triangle ABC}=\dfrac12\cdot AB^2=68$。}
  \end{solutionblock}
\end{problem}


\needspace{8\baselineskip}

\begin{problem}[points=10]
  已知 $A(1,2)$，$B(9,6)$。等腰 Rt$\triangle ABC$ 中 $\angle C = 90°$，$C$ 在第一象限且在 $AB$ 上方，$P$ 为 $AB$ 的中点。\\
(1) 求点 $C$ 的坐标；\\
(2) 求 $\angle APC$ 的度数和 $\triangle ABC$ 的面积。

  \answerarea[70mm]
  \setcounter{solstep}{0}
  \begin{solutionblock}
    \textbf{答案：}$C(3,\,8)$，$\angle APC=90°$，$S_{\triangle ABC}=20$\par\medskip
    \solstep{找斜边中点}{因为 $\angle C=90°$ 且 $AC=BC$，$AB$ 是斜边。$P$ 为 $AB$ 中点，所以 $P(5,4)$。}
    \solstep{求点 $C$}{$A$ 到 $P$ 的水平差为 $4$，竖直差为 $2$。过 $P$ 作与 $AB$ 垂直的方向，水平、竖直差互换为 $2$ 和 $4$。$C$ 在 $AB$ 上方，故 $C(5-2,4+4)=(3,8)$。}
    \solstep{求角度}{$PA^2=4^2+2^2=20$，$PC^2=2^2+4^2=20$，且一线三垂直说明 $PA\perp PC$，所以 $\angle APC=90°$。}
    \solstep{求面积}{$AC^2=(3-1)^2+(8-2)^2=40$，所以 $AC=BC=\sqrt{40}$。$S_{\triangle ABC}=\dfrac12\cdot AC\cdot BC=20$。}
  \end{solutionblock}
\end{problem}



\end{document}
